\documentclass[11pt, a4paper]{article}

\usepackage[utf8]{inputenc}
\usepackage[T1]{fontenc}
\usepackage[ngerman]{babel}
\usepackage{amsmath, amssymb, amsthm}
\usepackage{geometry}
\usepackage{pgfplots}
\pgfplotsset{compat=1.18}
\usepackage{booktabs}
\usepackage{array}
\usepackage{xcolor}
\usepackage{tcolorbox}
\tcbuselibrary{skins, breakable}
\usepackage{hyperref}
\usepackage{microtype}
\usepackage{parskip}
\usepackage{enumitem}

\geometry{left=2.5cm, right=2.5cm, top=2.5cm, bottom=2.5cm}

% Colors
\definecolor{fragecol}{RGB}{30, 90, 160}
\definecolor{antworttop}{RGB}{240, 246, 255}
\definecolor{antwortbottom}{RGB}{225, 237, 255}
\definecolor{sectioncol}{RGB}{20, 70, 130}

% Q&A boxes
\tcbset{
  frageboxstyle/.style={
    enhanced,
    colback=fragecol!10,
    colframe=fragecol,
    fonttitle=\bfseries\color{white},
    coltitle=white,
    attach boxed title to top left={yshift=-2mm, xshift=4mm},
    boxed title style={colback=fragecol, rounded corners},
    breakable,
    top=4mm,
  },
  antwortboxstyle/.style={
    enhanced,
    colback=antworttop,
    colframe=fragecol!40,
    breakable,
    top=2mm, bottom=2mm,
  }
}

\newenvironment{frage}[2]{%
  \begin{tcolorbox}[frageboxstyle, title={Frage}]
  \textbf{#1}\\[3pt]
  {\small\itshape\color{black!55}\(\triangleright\) #2}
  \end{tcolorbox}
  \begin{tcolorbox}[antwortboxstyle]
}{%
  \end{tcolorbox}
  \vspace{4mm}
}

\hypersetup{
  colorlinks=true,
  linkcolor=fragecol,
  urlcolor=fragecol,
}

\title{%
  \textbf{\color{sectioncol}CS229 -- Verständnisfragen}\\[4pt]
  \large Strukturiert nach Thematik
}
\author{}
\date{Stand: \today}

\begin{document}
\maketitle
\tableofcontents
\newpage

% ============================================================
\section{Logistic Regression -- Loss-Funktion und Gradient}
% ============================================================

\begin{frage}{Warum gilt die Umrechnung zwischen Gradient und Loss-Funktion?}{PS2-1 »Training Stability«, Notebook-Abschnitt zur Herleitung der Loss-Funktion aus dem Gradienten}

Der Gradient der logistischen Regression lautet:
\[
  \nabla_\theta J(\theta)
  = -\frac{1}{m} \sum_{i=1}^{m}
      \frac{y^{(i)}\, x^{(i)}}{1 + \exp\!\bigl(y^{(i)}\,\theta^\top x^{(i)}\bigr)}
\]

\textbf{Behauptung:} Dieser Gradient gehört zur Loss-Funktion
\[
  \ell(\theta)
  = \frac{1}{m} \sum_{i=1}^{m} \log\!\bigl(1 + \exp\!\bigl(-y^{(i)}\,\theta^\top x^{(i)}\bigr)\bigr).
\]

\textbf{Nachrechnen durch Ableiten:}

\textit{Schritt 1 -- Einen Term ableiten:}
\[
  \frac{\partial}{\partial\theta}
    \log\!\bigl(1 + \exp(-y\,\theta^\top x)\bigr)
  = \frac{1}{1 + e^{-y\,\theta^\top x}}
    \cdot e^{-y\,\theta^\top x} \cdot (-y\,x)
\]

\textit{Schritt 2 -- Vorfaktor vereinfachen} (mit $a := y\,\theta^\top x$):
\[
  \frac{e^{-a}}{1 + e^{-a}}
  = \frac{e^{-a}}{1 + e^{-a}} \cdot \frac{e^{a}}{e^{a}}
  = \frac{1}{e^{a} + 1}
  = \frac{1}{1 + e^{a}}
\]

\textit{Schritt 3 -- Einsetzen:}
\[
  \frac{\partial}{\partial\theta}
    \log\!\bigl(1 + e^{-y\,\theta^\top x}\bigr)
  = \frac{-y\,x}{1 + e^{y\,\theta^\top x}}
\]

\textit{Schritt 4 -- Gesamt:}
\[
  \nabla_\theta\,\ell(\theta)
  = \frac{1}{m} \sum_{i=1}^{m}
      \frac{-y^{(i)}\,x^{(i)}}{1 + \exp\!\bigl(y^{(i)}\,\theta^\top x^{(i)}\bigr)}
  = -\frac{1}{m} \sum_{i=1}^{m}
      \frac{y^{(i)}\,x^{(i)}}{1 + \exp\!\bigl(y^{(i)}\,\theta^\top x^{(i)}\bigr)}
\]

Das stimmt mit dem angegebenen Gradienten überein. Die entscheidende Vereinfachung ist
$e^{-a}/(1+e^{-a}) = 1/(1+e^{a})$.

\end{frage}

\begin{frage}{Was bedeutet \(\ell\) (``ell'')?}{Allgemeine CS229-Notation; tritt in PS2-1 Notebook und allen Loss-Herleitungen auf}

$\ell$ ist der Name für die \textbf{Loss-Funktion}, abgeleitet vom Buchstaben \textit{l} (kleines L,
kursiv), der in der Mathematik häufig für ``Loss'' oder ``Likelihood'' steht.

In diesem Kontext bezeichnet $\ell(\theta)$ die \textbf{negative Log-Likelihood} der logistischen
Regression:
\[
  \ell(\theta)
  = \frac{1}{m} \sum_{i=1}^{m}
    \log\!\bigl(1 + \exp\!\bigl(-y^{(i)}\,\theta^\top x^{(i)}\bigr)\bigr),
\]
die minimiert werden soll.

\end{frage}

% ============================================================
\section{Lineare Trennbarkeit und Konvergenz}
% ============================================================

\begin{frage}{Warum gilt: Dataset ist linear trennbar $\Leftrightarrow$ $y^{(i)}\,\theta^\top x^{(i)} > 0$ für alle $i$?}{PS2-1 »Training Stability«, Teil (b) -- Ursache der Divergenz; Grundlage in main\_notes.pdf, Kap.\,6.3 (Funktionaler Margin, S.\,61)}

Die Entscheidungsgrenze ist die Hyperebene $\{\,x : \theta^\top x = 0\,\}$. Ein Punkt liegt auf
einer der zwei Seiten je nach Vorzeichen von $\theta^\top x$:

\begin{center}
\begin{tabular}{cccl}
  \toprule
  $y^{(i)}$ & $\theta^\top x^{(i)}$ & $y^{(i)}\cdot\theta^\top x^{(i)}$ & Klassifikation \\
  \midrule
  $+1$ & $> 0$ & $> 0$ & korrekt \\
  $-1$ & $< 0$ & $> 0$ & korrekt \\
  $+1$ & $< 0$ & $< 0$ & falsch \\
  $-1$ & $> 0$ & $< 0$ & falsch \\
  \bottomrule
\end{tabular}
\end{center}

$y^{(i)}\,\theta^\top x^{(i)} > 0$ bedeutet also genau: Punkt $i$ liegt auf der richtigen Seite der
Trennhyperebene. Gilt das für alle $i$, trennt $\theta$ alle Punkte korrekt -- das ist die Definition
von \textbf{linearer Trennbarkeit}.

\textbf{Im Skript:} \texttt{main\_notes.pdf}, Kap.\,6.3, S.\,61:
\begin{quote}
\itshape
``if $y^{(i)}(w^\top x + b) > 0$, then our prediction on this example is correct.''
\end{quote}
Dies ist der sogenannte \textbf{funktionale Margin} $\hat{\gamma}^{(i)} = y^{(i)}(w^\top x + b)$.

\end{frage}

\begin{frage}{Warum konvergiert Gradient Descent nicht für linear trennbare Datasets?}{PS2-1 »Training Stability«, Teil (b) -- Erklärung des Divergenzverhaltens von Dataset~B}

\textbf{Schritt 1: Die Loss-Funktion hat kein Minimum.}

Sei das Dataset trennbar, d.\,h.\ es gibt $\theta^*$ mit $y^{(i)}\,\theta^{*\top} x^{(i)} > 0$
für alle $i$. Betrachte $c\,\theta^*$ für wachsendes $c > 0$:
\[
  \ell(c\,\theta^*)
  = \frac{1}{m} \sum_{i=1}^{m}
    \underbrace{\log\!\bigl(1 + \exp\!\bigl(-c\cdot\underbrace{y^{(i)}\,\theta^{*\top} x^{(i)}}_{>0}\bigr)\bigr)}_{\to\, 0 \text{ für } c\to\infty}
\]

Da der Exponent $-c\cdot(\text{positiv})$ immer negativer wird:
\begin{align*}
  c = 1:   & \quad \log(1 + e^{-1}) \approx 0{,}31 \\
  c = 10:  & \quad \log(1 + e^{-10}) \approx 0{,}000045 \\
  c \to \infty: & \quad \to 0, \text{ aber nie } = 0
\end{align*}

Der Loss sinkt unbegrenzt, erreicht aber sein Infimum $0$ nie. Es gibt also \textbf{kein echtes
Minimum} -- der Gradient wird nie exakt $\mathbf{0}$.

\textbf{Schritt 2: Konsequenz für Gradient Descent.}

Da kein endliches Minimum existiert, gilt $\|\theta\| \to \infty$. Das Abbruchkriterium
$\|\theta_{t+1} - \theta_t\| < \varepsilon$ wird nie erfüllt.

\textbf{Schritt 3: Nicht-trennbarer Fall zum Vergleich.}

Wenn es mindestens einen falsch klassifizierten Punkt gibt, hilft Skalieren nicht: der Loss dieses
Punktes wird durch größeres $c$ \emph{größer}, nicht kleiner. Der Loss hat daher ein echtes Minimum
bei endlichem $\theta$, und Gradient Descent konvergiert dorthin.

\end{frage}

\begin{frage}{Ist das Divergenzproblem generell bei Logistic Regression?}{Folge aus PS2-1 »Training Stability«, Teil (c) -- Gegenmassnahmen; allgemeine Frage zur Methode}

Ja, aber nur wenn das Dataset \textbf{linear trennbar} ist. Das ist kein Implementierungsfehler,
sondern ein mathematisches Problem: die Log-Likelihood hat kein endliches Maximum.

\textbf{Lösungen in der Praxis:}

\begin{enumerate}[leftmargin=*, label=(\arabic*)]
  \item \textbf{L2-Regularisierung} (häufigstes Vorgehen):
    \[
      \ell_{\text{reg}}(\theta)
      = \ell(\theta) + \lambda\|\theta\|^2
    \]
    Für $\|\theta\| \to \infty$ gilt $\lambda\|\theta\|^2 \to \infty$, also gibt es ein endliches
    Minimum.

  \item \textbf{Early Stopping} -- Training abbrechen bevor $\theta$ zu groß wird.

  \item \textbf{Abnehmende Lernrate} -- Schritte werden mit der Zeit kleiner.
\end{enumerate}

In \texttt{scikit-learn} ist L2-Regularisierung standardmäßig aktiv (\texttt{C=1.0}), weshalb das
Problem dort normalerweise nicht auftritt.

\end{frage}

% ============================================================
\section{SVM und Hinge Loss}
% ============================================================

\begin{frage}{Was ist Hinge Loss?}{PS2-1 »Training Stability«, Teil (d) -- Vergleich LR mit SVM als Alternative bei linear trennbaren Daten; Theorie in main\_notes.pdf, Kap.\,6}

Hinge Loss ist die Loss-Funktion der SVMs. Definition für einen Punkt:
\[
  L\!\bigl(\theta, x^{(i)}, y^{(i)}\bigr)
  = \max\!\bigl(0,\; 1 - y^{(i)}\,\theta^\top x^{(i)}\bigr)
\]

\textbf{Verhalten:}
\begin{center}
\begin{tabular}{lll}
  \toprule
  $y\,\theta^\top x$ & Hinge Loss & Interpretation \\
  \midrule
  $\geq 1$ & $0$ & korrekt, mit ausreichend Abstand (Margin $\geq 1$) \\
  $= 0$    & $1$ & Punkt liegt auf der Entscheidungsgrenze \\
  $= -1$   & $2$ & falsch klassifiziert \\
  \bottomrule
\end{tabular}
\end{center}

\begin{center}
\begin{tikzpicture}
\begin{axis}[
  width=10cm, height=6cm,
  xlabel={$z = y\,\theta^\top x$},
  ylabel={Loss},
  xmin=-2.2, xmax=2.5,
  ymin=-0.1, ymax=2.8,
  grid=both,
  grid style={line width=0.3pt, draw=gray!30},
  major grid style={line width=0.5pt, draw=gray!60},
  legend pos=north east,
  legend style={font=\small},
  tick label style={font=\small},
  label style={font=\small},
  axis lines=left,
]
  % Hinge Loss: max(0, 1-z)
  \addplot[blue, very thick, domain=-2.2:1] {1 - x};
  \addplot[blue, very thick, domain=1:2.5]  {0};
  \addlegendentry{Hinge Loss: $\max(0, 1-z)$}

  % Logistic Loss: log(1 + exp(-z)) / log(2)  [scaled for comparison]
  \addplot[red, dashed, thick, domain=-2.2:2.5, samples=100]
    {ln(1 + exp(-x)) / ln(2)};
  \addlegendentry{Log. Loss: $\log_2(1+e^{-z})$}

  % Mark z=1
  \draw[dotted, gray] (axis cs:1,0) -- (axis cs:1,1);
  \node[below, font=\small] at (axis cs:1,0) {$1$};
\end{axis}
\end{tikzpicture}
\end{center}

\textbf{Vergleich mit logistischem Loss:}
\begin{align*}
  \text{Logistic Loss:} &\quad \log(1 + e^{-z}) > 0 \text{ immer} \\
  \text{Hinge Loss:}    &\quad \max(0, 1-z) = 0 \text{ für } z \geq 1
\end{align*}

Deshalb hat das SVM das Konvergenzproblem bei trennbaren Datasets \emph{nicht}: Sobald alle Punkte
$y^{(i)}\,\theta^\top x^{(i)} \geq 1$ erfüllen, ist der Loss exakt $0$ -- kein weiterer Gradient,
kein Wachstum von $\theta$.

\end{frage}

% ============================================================
\section{Mathematische Notation}
% ============================================================

\begin{frage}{Was bedeutet $\bigl|\bigl\{\,i \in I_{a,b}\,\bigr\}\bigr|$?}{PS2-2 »Model Calibration«, Kalibrierungsdefinition -- $I_{a,b}$ bezeichnet den Bucket von Beispielen mit Vorhersage in $[a,b]$}

Das ist die \textbf{Kardinalität} (Anzahl der Elemente) einer Menge:
\[
  \bigl|\bigl\{\,i \in I_{a,b}\,\bigr\}\bigr|
  \;=\; \text{Anzahl der Indizes } i \text{, die in der Menge } I_{a,b} \text{ enthalten sind.}
\]

\textbf{Beispiel:} Wenn $I_{a,b} = \{2, 5, 7, 9\}$, dann $|I_{a,b}| = 4$.

Allgemein: $|S|$ bedeutet ``wie viele Elemente hat Menge $S$''.

\end{frage}

\begin{frage}{Was bedeutet $\mathbb{I}\bigl\{y^{(i)} = 1\bigr\}$?}{PS2-2 »Model Calibration«, Kalibrierungsformel -- allgemeine CS229-Notation für Indikatorfunktionen}

Das ist die \textbf{Indikatorfunktion}:
\[
  \mathbb{I}\bigl\{y^{(i)} = 1\bigr\}
  = \begin{cases}
      1 & \text{falls } y^{(i)} = 1 \\
      0 & \text{sonst}
    \end{cases}
\]

Sie gibt $1$ zurück wenn die Bedingung wahr ist, sonst $0$.

\textbf{Beispiel} mit $m=4$ Datenpunkten, $y = (1, 0, 1, 1)$:
\begin{align*}
  \mathbb{I}\{y^{(1)} = 1\} &= 1 \\
  \mathbb{I}\{y^{(2)} = 1\} &= 0 \\
  \mathbb{I}\{y^{(3)} = 1\} &= 1 \\
  \mathbb{I}\{y^{(4)} = 1\} &= 1
\end{align*}

In Python entspricht das \texttt{(y == 1).astype(int)} bzw.\ direkt \texttt{y == 1} in einem
booleschen Ausdruck.

\end{frage}

% ============================================================
\section{Model Calibration (PS2-2)}
% ============================================================

\begin{frage}{Impliziert perfekte Kalibrierung perfekte Genauigkeit -- und umgekehrt?}{PS2-2 »Model Calibration«, Teil (b) -- Beweis der Unabhängigkeit beider Eigenschaften}

\textbf{Antwort: Nein, in beide Richtungen.}

\medskip
\textbf{Richtung 1: Kalibrierung impliziert \emph{nicht} Genauigkeit.}

\textit{Gegenbeispiel:} Alle $y^{(i)} = 1$ (nur positive Klasse).

Für den Bucket $(a, b) = (0{,}5,\; 1)$ muss bei perfekter Kalibrierung gelten:
\[
  \frac{\sum_{i \in I_{a,b}} \mathbb{I}\{y^{(i)}=1\}}{|\{i \in I_{a,b}\}|}
  = \frac{\sum_{i \in I_{a,b}} P(y^{(i)}=1 \mid x^{(i)};\,\theta)}{|\{i \in I_{a,b}\}|}
\]

Die linke Seite ist $1$ (alle Punkte haben $y=1$). Da $h(x) = \sigma(\theta^\top x) < 1$ immer gilt,
muss die rechte Seite $< 1$ sein -- das Modell sagt für manche Punkte $h(x) < 0{,}5$ vorher: also
\textbf{falsch klassifiziert}.

\medskip
\textbf{Richtung 2: Genauigkeit impliziert \emph{nicht} Kalibrierung.}

Perfekte Genauigkeit bedeutet alle Punkte korrekt klassifiziert, also Anteil $y=1$ im Bucket $= 1$.
Aber das Modell könnte trotzdem $h(x) = 0{,}6$ vorhersagen:
\[
  \underbrace{\frac{\sum_{i \in I_{a,b}} \mathbb{I}\{y^{(i)}=1\}}{|\{i \in I_{a,b}\}|}}_{= 1}
  \neq
  \underbrace{\frac{\sum_{i \in I_{a,b}} P(y^{(i)}=1 \mid x^{(i)};\,\theta)}{|\{i \in I_{a,b}\}|}}_{< 1}
\]

Kalibrierung verletzt, obwohl alle Entscheidungen korrekt sind.

\medskip
\textbf{Kernaussage:} Kalibrierung prüft ob Wahrscheinlichkeiten realistisch sind. Genauigkeit prüft
nur ob die Entscheidung ($> 0{,}5$ oder nicht) stimmt. Beides ist voneinander unabhängig.

\end{frage}

% ============================================================
\section{Kursstruktur}
% ============================================================

\begin{frage}{Greifen die PS2-Aufgaben auf Wissen aus späteren Kapiteln vor?}{Übergreifende Frage zur Kursstruktur -- entstanden beim Durcharbeiten aller PS2-Aufgaben}

Ja. Mehrere PS2-Aufgaben setzen Konzepte voraus, die im Kursverlauf erst später behandelt werden:

\begin{center}
\begin{tabular}{p{2.5cm} p{3.5cm} p{4cm} p{4cm}}
  \toprule
  \textbf{Problem} & \textbf{Erwartete Notes} & \textbf{Tatsächlich benötigt} & \textbf{Hinweis} \\
  \midrule
  PS2-1 & Kap.\,2 (LR) & Kap.\,6 (Hinge Loss/SVM) für Teil~(d); Kap.\,9 (Regularisierung) für Teil~(c) & Teil~(d) setzt SVM-Konzepte voraus \\
  \addlinespace
  PS2-2 & Kap.\,2 (LR, MLE) & -- & vollständig in Kap.\,2 \\
  \addlinespace
  PS2-3 & Kap.\,2/9 & \textbf{Kap.\,9.4} (Bayesian Statistics) als Hauptquelle & Kap.\,9.4 kommt später im Kurs \\
  \addlinespace
  PS2-4 & -- & Kap.\,5+6 (Kernels, Mercer) vollständig & komplett aus Kap.\,5+6 \\
  \addlinespace
  PS2-5 & -- & Kap.\,5+6 (Kernel Trick, Representer Theorem) & komplett aus Kap.\,5+6 \\
  \addlinespace
  PS2-6 & Kap.\,4.2 (NB) & Kap.\,6 (SVM) für Teil~(d) & Hybrid-Problem \\
  \bottomrule
\end{tabular}
\end{center}

\textbf{Fazit:} PS2-4 und PS2-5 sind komplett aus Kap.\,5+6 (SVMs/Kernels), das nach PS1 kommt. PS2-3
erfordert Kap.\,9.4 (Bayesian Statistics) als primäre Quelle -- ein deutlicher Vorgriff auf spätere
Kursinhalte.

\end{frage}

\end{document}
